\chapter{Persistencia extendida y superficies de mundo de cuerdas y branas}
\label{mdp4:chap:teoria-m}

\section{Persistencia extendida y superficie de mundo}

Una cuerda física es un objeto unidimensional extendido cuya evolución barre
una superficie de mundo. Cada punto de la cuerda posee una coordenada interna
y cada instante añade una coordenada temporal; la combinación de ambas forma
el dominio sobre el que viven la métrica inducida, la tensión y los campos de
fondo. Una cuerda cerrada tiene sección espacial circular. Una cuerda abierta
tiene extremos sometidos a condiciones de frontera. Sus vibraciones organizan
estados y su tensión convierte el área recorrida por la superficie de mundo en
acción.

Una brana de dimensión \(p\) es un objeto extendido cuya historia forma un
volumen de mundo de dimensión \(p+1\). Las branas reciben extremos de cuerdas,
transportan campos sobre su volumen y pueden envolver ciclos de un espacio
compacto. Esta definición fija tres datos diferentes: la dimensión del objeto,
la dimensión de su historia y las condiciones que gobiernan su frontera. La
tensión, las cargas y el espectro pertenecen a la realización dinámica que
completa esos datos geométricos.

HMT llega a esta frontera desde una arquitectura de rutas persistentes. Una
ruta conserva orientación, hoja, fase, cociente, acarreo y memoria de los
cruces que la han prolongado. El paso hacia un objeto extendido asigna esa
genealogía a una sección distribuida sobre una coordenada interna. La
composición de rutas se realiza entonces como costura de segmentos; la memoria
de frontera registra enrollamientos e intersecciones; la acción mide la
persistencia completa sobre la superficie recorrida. Este mapa conserva el
orden causal: primero se construye el estado enriquecido y después se declara
su realización como cuerda, brana o campo.

La acción de Polyakov ofrece una forma precisa para el codominio físico:
\[
 S_{\rm P}[X,h]
 =-\frac{T}{2}\int d^2\sigma\,
 \sqrt{-h}\,h^{ab}g_{MN}(X)
 \partial_aX^M\partial_bX^N.
\]
Aquí \(T\) es la tensión, \(h_{ab}\) la métrica de la superficie de mundo y
\(X^M\) el mapa hacia el espacio objetivo. Una realización HMT de cuerda debe
producir estos objetos, su espacio de estados y sus simetrías desde un mapa
tipado; la estructura de ruta suministra el dominio genealógico sobre el que
ese mapa actúa.

\section{Pantallas exactas en dimensiones \(11\), \(10\), \(26\) y \(24\)}

La primera pantalla parte del código ternario extendido
\(C_W=[12,6,6]_3\). La perforación de una coordenada produce el código
\(G_{11}=[11,6,5]_3\). Sus bolas de Hamming de radio dos poseen cardinal
\[
 V_3(11,2)=1+22+220=243=3^5,
\]
y las \(3^6\) palabras del código satisfacen
\[
 3^6V_3(11,2)=3^{11}.
\]
Por tanto, esas bolas particionan exactamente el espacio ambiente. El entero
once mide aquí la longitud de un código perfecto y conserva su tipo
combinatorio.

Una segunda pantalla usa el módulo real de permutación \(V_{12}\) asociado a
los doce vértices icosaédricos. El complemento del vector uniforme se
descompone como
\[
 V_{11}=V_3\oplus V_{3'}\oplus V_5.
\]
El estado dodecafásico \(K\) polariza una recta \(\ell_K\subset V_3\) y
determina el proyector
\[
 P_{10}=P_{11}-\frac{u_Ku_K^{\mathsf T}}{\lVert u_K\rVert^2}.
\]
De él descienden las descomposiciones exactas
\[
 V_{11}=\ell_K\oplus V_{10},\qquad
 V_{10}=V_{3'}\oplus W_7(K),
\]
con dimensiones \(11=1+3+7\) y \(10=3+7\). El paso
\(12\to11\) elimina el modo uniforme; el paso \(11\to10\) elimina la
dirección marcada por \(K\). Una compactificación física añade una recta
temporal, una retícula estable de rango seis y el transporte que preserve las
tres direcciones extendidas.

La tercera pantalla parte del retículo
\[
 L_{26}=\Lambda_{24}\oplus U\cong II_{25,1},
\]
donde \(U\) es el plano hiperbólico entero. Para un vector isotrópico
primitivo \(e_+\), la reducción
\[
 e_+^\perp/\mathbb Ze_+\cong\Lambda_{24}
\]
es exacta. La condición de ortogonalidad retira una dirección y el cociente
retira la dirección nula contenida en su ortogonal; así se obtiene el rango
veinticuatro. Código, módulo real y retículo pertenecen a categorías
distintas. Sus mapas explícitos permiten relacionar sus rangos conservando la
procedencia de cada uno.

La incidencia entre las seis caras y los ocho vértices del cubo aporta además
una pantalla de firma espacial y temporal. Si \(M\) es su matriz de
incidencia, entonces
\[
 MM^{\mathsf T}=12P_u+4P_-.
\]
El cociente por el núcleo posee rango cuatro y se descompone como
\(1\oplus3\). Al declarar la recta uniforme como dirección temporal, la forma
normalizada queda \(\operatorname{diag}(-1,1,1,1)\). La incidencia demuestra
el rango y la razón espectral; la orientación declarada fija la firma.

\section{Momento, enrollamiento y radio recíproco}

En una coordenada compacta de radio \(R\), el entero \(m\) registra momento y
el entero \(w\) registra enrollamiento. La forma
\[
 E_R(m,w)=\frac{m^2}{R^2}+w^2R^2
\]
posee la involución
\[
 \mathsf T(m,w,R)=(w,m,R^{-1}),
 \qquad
 E_R(m,w)=E_{R^{-1}}(w,m).
\]
En \(\ell^2(\mathbb Z^2)\), el intercambio
\(U_T\lvert m,w\rangle=\lvert w,m\rangle\) proporciona la equivalencia
unitaria
\[
 U_TH(R)U_T^{-1}=H(R^{-1}).
\]
Esta es la residencia algebraica de la dualidad de radio: momento y
enrollamiento son dos lecturas recíprocas de un mismo sector reticular. La
realización como dualidad T incorpora la escala de cuerda, los osciladores y
la acción sobre el espacio físico de estados. En el radio autodual ambas
lecturas comparten la misma escala y el intercambio actúa dentro del mismo
espectro.

\section{Interfaz en dimensión \(11\): terna espectral, supersimetría y operador \(\mathcal Q^2=H\)}

La teoría M organiza un codominio físico de once dimensiones con métrica
\(g_{MN}\), tres-forma \(C_{MNP}\), gravitino \(\psi_M\), acción y
transformaciones de supersimetría. El puente se escribe como un mapa
\[
 \mathfrak R_M:\mathcal X_{\rm HMT}^{\rm enr}
 \longrightarrow\mathcal X_{11}^{\rm phys}.
\]
Su fuente conserva fase, orientación, costura y memoria; su codominio añade
los campos y la dinámica undecadimensional. La compactificación de una
dirección debe transportar esos campos a una pantalla de diez dimensiones,
preservar el límite de baja energía y entrelazar las dualidades declaradas.

Una compactificación candidata queda tipada por una terna
\((x_\infty,\kappa,\mathcal U)\). La sección \(x_\infty\) representa una
persistencia prolongable; \(\kappa\) la lleva a un espacio de módulos;
\(\mathcal U\) transporta métrica, orientación, registro y frontera. Su
admisibilidad exige prolongación a toda profundidad, conmutación con la
dualidad, descenso del espectro al cociente y conservación de dominios y
codominios de los observables. Estas condiciones convierten una analogía de
rangos en un problema matemático verificable.

Las cuerdas aparecen tras compactificar el sector undecadimensional; las
branas aparecen como realizaciones extendidas acopladas a la tres-forma y a
sus duales. La rama excepcional que llega a \(\Lambda_{24}\),
\(V^\natural\) y el Monstruo conserva la misma fuente enriquecida mediante
otro funtor. La interfaz dimensional y la excepcional son prolongaciones
diferentes de una genealogía común, cada una con su propio codominio.

\begin{resultado}
La perfección de \(G_{11}\), la polarización \(11\to10\), la reducción
isotrópica \(26\to24\), la pantalla cúbica de firma \((3,1)\) y la
involución de radio son resultados algebraicos demostrados. Juntos forman la
residencia exacta de rangos, cocientes y dualidades sobre la que se construye
la interfaz hacia cuerdas, branas y teoría M.
\end{resultado}

\begin{lecturafisica}
Una historia extendida conserva posición, enrollamiento, orientación y
frontera dentro de una misma persistencia. La compactificación determina qué
direcciones aparecen como extensión macroscópica y cuáles actúan como grados
internos. La cuerda lee esa persistencia sobre una superficie de mundo; la
brana la lee sobre un volumen de mundo; la teoría M coordina ambas lecturas
mediante campos undecadimensionales.
\end{lecturafisica}

\begin{frontera}
La realización se refuta si el mapa \(\mathfrak R_M\) altera los tipos de las
pantallas, si la dualidad deja de actuar unitariamente sobre el espectro, si
la compactificación pierde prolongación o si los campos incumplen acción y
supersimetría. El cierre físico requiere tensión, osciladores, amplitudes,
anomalías controladas y un límite gravitatorio coherente para cada
compactificación declarada.
\end{frontera}
